% Estilos comunes de los diagramas (TikZ puro, sin tikz-uml).
% Se carga con % Estilos comunes de los diagramas (TikZ puro, sin tikz-uml).
% Se carga con % Estilos comunes de los diagramas (TikZ puro, sin tikz-uml).
% Se carga con % Estilos comunes de los diagramas (TikZ puro, sin tikz-uml).
% Se carga con \input{estilos_uml} desde cada diagrama standalone.
\usepackage[T1]{fontenc}
\usepackage[utf8]{inputenc}
\usepackage{lmodern}
\usepackage{underscore}   % permite escribir identificadores con "_" sin escapar
\usepackage{xcolor}
\DeclareUnicodeCharacter{2192}{\ensuremath{\rightarrow}}
\DeclareUnicodeCharacter{2190}{\ensuremath{\leftarrow}}
\DeclareUnicodeCharacter{2194}{\ensuremath{\leftrightarrow}}
\DeclareUnicodeCharacter{2265}{\ensuremath{\geq}}
\DeclareUnicodeCharacter{2264}{\ensuremath{\leq}}
\DeclareUnicodeCharacter{2260}{\ensuremath{\neq}}
\DeclareUnicodeCharacter{00D7}{\ensuremath{\times}}
\DeclareUnicodeCharacter{21D2}{\ensuremath{\Rightarrow}}
\DeclareUnicodeCharacter{2026}{\ldots}
\usetikzlibrary{shapes.multipart,shapes.geometric,shapes.misc,arrows.meta,positioning,fit,backgrounds,calc}

% Paleta (sobria, imprime bien en B/N)
\definecolor{cPep}{HTML}{E3EEF9}     % agente / PEP
\definecolor{cPdp}{HTML}{FCEBD9}     % API / PDP
\definecolor{cDato}{HTML}{E4F2E2}    % datos / persistencia
\definecolor{cExt}{HTML}{EEE6F5}     % servicios externos
\definecolor{cUi}{HTML}{FFF6CC}      % frontend
\definecolor{cOff}{HTML}{EDEDED}     % offline
\definecolor{cBorde}{HTML}{3A3A3A}

\tikzset{
  % --- UML: clase de 3 compartimentos
  clase/.style={rectangle split, rectangle split parts=3, draw=cBorde, fill=white,
    font=\ttfamily\scriptsize, inner sep=2.5pt, text width=#1,
    rectangle split part align={center,left,left}},
  clase/.default=6cm,
  % --- UML: objeto (2 compartimentos, nombre subrayado)
  objeto/.style={rectangle split, rectangle split parts=2, draw=cBorde, fill=white,
    font=\ttfamily\scriptsize, inner sep=2.5pt, text width=#1,
    rectangle split part align={center,left}},
  objeto/.default=5cm,
  % --- relaciones UML
  herencia/.style={-{Triangle[open,length=7pt,width=8pt]}, draw=cBorde, thick},
  realiza/.style={-{Triangle[open,length=7pt,width=8pt]}, draw=cBorde, thick, dashed},
  compone/.style={{Diamond[length=9pt,width=6pt]}-{Stealth[length=5pt]}, draw=cBorde, thick},
  agrega/.style={{Diamond[open,length=9pt,width=6pt]}-{Stealth[length=5pt]}, draw=cBorde, thick},
  usa/.style={-{Stealth[length=5pt]}, draw=cBorde, dashed},
  asocia/.style={-{Stealth[length=5pt]}, draw=cBorde},
  enlace/.style={draw=cBorde, thick},
  rotulo/.style={font=\ttfamily\tiny, fill=white, inner sep=1pt},
  % --- componentes / flujo
  comp/.style={draw=cBorde, rounded corners=2pt, align=center, font=\small, inner sep=4pt, minimum height=9mm},
  paso/.style={draw=cBorde, rounded corners=3pt, align=center, font=\footnotesize, inner sep=3pt, text width=#1},
  paso/.default=3.4cm,
  decision/.style={draw=cBorde, diamond, aspect=2.2, align=center, font=\footnotesize, inner sep=1pt},
  almacen/.style={draw=cBorde, cylinder, shape border rotate=90, aspect=0.18, align=center,
    font=\footnotesize, minimum width=2.6cm, minimum height=1.2cm},
  zona/.style={draw=cBorde, dashed, rounded corners=4pt, inner sep=8pt},
  zonatit/.style={font=\bfseries\small, anchor=north west, inner sep=2pt},
  flujo/.style={-{Stealth[length=6pt]}, draw=cBorde, thick},
  flujofallo/.style={-{Stealth[length=6pt]}, draw=red!60!black, thick, dashed},
  proto/.style={font=\scriptsize\ttfamily, fill=white, inner sep=1pt, align=center},
}

% \umlclase[opciones]{nombre-nodo}{x,y}{estereotipo}{Nombre}{atributos}{operaciones}
% atributos/operaciones: líneas separadas con \\ ; usar \vacio si no hay.
\newcommand{\vacio}{\phantom{x}}
\newcommand{\umlclase}[7][]{%
  \node[clase,#1] (#2) at (#3) {\makebox[\linewidth][c]{%
    \begin{tabular}{@{}c@{}}\if\relax\detokenize{#4}\relax\else\textit{«#4»}\\\fi\textbf{\normalsize #5}\end{tabular}}%
    \nodepart{two}\begin{tabular}{@{}l@{}}#6\end{tabular}%
    \nodepart{three}\begin{tabular}{@{}l@{}}#7\end{tabular}};}

% \umlobjeto[opciones]{nombre-nodo}{x,y}{nombre:Clase}{slots}
\newcommand{\umlobjeto}[5][]{%
  \node[objeto,#1] (#2) at (#3) {\makebox[\linewidth][c]{\underline{\textbf{#4}}}%
    \nodepart{two}\begin{tabular}{@{}l@{}}#5\end{tabular}};}
 desde cada diagrama standalone.
\usepackage[T1]{fontenc}
\usepackage[utf8]{inputenc}
\usepackage{lmodern}
\usepackage{underscore}   % permite escribir identificadores con "_" sin escapar
\usepackage{xcolor}
\DeclareUnicodeCharacter{2192}{\ensuremath{\rightarrow}}
\DeclareUnicodeCharacter{2190}{\ensuremath{\leftarrow}}
\DeclareUnicodeCharacter{2194}{\ensuremath{\leftrightarrow}}
\DeclareUnicodeCharacter{2265}{\ensuremath{\geq}}
\DeclareUnicodeCharacter{2264}{\ensuremath{\leq}}
\DeclareUnicodeCharacter{2260}{\ensuremath{\neq}}
\DeclareUnicodeCharacter{00D7}{\ensuremath{\times}}
\DeclareUnicodeCharacter{21D2}{\ensuremath{\Rightarrow}}
\DeclareUnicodeCharacter{2026}{\ldots}
\usetikzlibrary{shapes.multipart,shapes.geometric,shapes.misc,arrows.meta,positioning,fit,backgrounds,calc}

% Paleta (sobria, imprime bien en B/N)
\definecolor{cPep}{HTML}{E3EEF9}     % agente / PEP
\definecolor{cPdp}{HTML}{FCEBD9}     % API / PDP
\definecolor{cDato}{HTML}{E4F2E2}    % datos / persistencia
\definecolor{cExt}{HTML}{EEE6F5}     % servicios externos
\definecolor{cUi}{HTML}{FFF6CC}      % frontend
\definecolor{cOff}{HTML}{EDEDED}     % offline
\definecolor{cBorde}{HTML}{3A3A3A}

\tikzset{
  % --- UML: clase de 3 compartimentos
  clase/.style={rectangle split, rectangle split parts=3, draw=cBorde, fill=white,
    font=\ttfamily\scriptsize, inner sep=2.5pt, text width=#1,
    rectangle split part align={center,left,left}},
  clase/.default=6cm,
  % --- UML: objeto (2 compartimentos, nombre subrayado)
  objeto/.style={rectangle split, rectangle split parts=2, draw=cBorde, fill=white,
    font=\ttfamily\scriptsize, inner sep=2.5pt, text width=#1,
    rectangle split part align={center,left}},
  objeto/.default=5cm,
  % --- relaciones UML
  herencia/.style={-{Triangle[open,length=7pt,width=8pt]}, draw=cBorde, thick},
  realiza/.style={-{Triangle[open,length=7pt,width=8pt]}, draw=cBorde, thick, dashed},
  compone/.style={{Diamond[length=9pt,width=6pt]}-{Stealth[length=5pt]}, draw=cBorde, thick},
  agrega/.style={{Diamond[open,length=9pt,width=6pt]}-{Stealth[length=5pt]}, draw=cBorde, thick},
  usa/.style={-{Stealth[length=5pt]}, draw=cBorde, dashed},
  asocia/.style={-{Stealth[length=5pt]}, draw=cBorde},
  enlace/.style={draw=cBorde, thick},
  rotulo/.style={font=\ttfamily\tiny, fill=white, inner sep=1pt},
  % --- componentes / flujo
  comp/.style={draw=cBorde, rounded corners=2pt, align=center, font=\small, inner sep=4pt, minimum height=9mm},
  paso/.style={draw=cBorde, rounded corners=3pt, align=center, font=\footnotesize, inner sep=3pt, text width=#1},
  paso/.default=3.4cm,
  decision/.style={draw=cBorde, diamond, aspect=2.2, align=center, font=\footnotesize, inner sep=1pt},
  almacen/.style={draw=cBorde, cylinder, shape border rotate=90, aspect=0.18, align=center,
    font=\footnotesize, minimum width=2.6cm, minimum height=1.2cm},
  zona/.style={draw=cBorde, dashed, rounded corners=4pt, inner sep=8pt},
  zonatit/.style={font=\bfseries\small, anchor=north west, inner sep=2pt},
  flujo/.style={-{Stealth[length=6pt]}, draw=cBorde, thick},
  flujofallo/.style={-{Stealth[length=6pt]}, draw=red!60!black, thick, dashed},
  proto/.style={font=\scriptsize\ttfamily, fill=white, inner sep=1pt, align=center},
}

% \umlclase[opciones]{nombre-nodo}{x,y}{estereotipo}{Nombre}{atributos}{operaciones}
% atributos/operaciones: líneas separadas con \\ ; usar \vacio si no hay.
\newcommand{\vacio}{\phantom{x}}
\newcommand{\umlclase}[7][]{%
  \node[clase,#1] (#2) at (#3) {\makebox[\linewidth][c]{%
    \begin{tabular}{@{}c@{}}\if\relax\detokenize{#4}\relax\else\textit{«#4»}\\\fi\textbf{\normalsize #5}\end{tabular}}%
    \nodepart{two}\begin{tabular}{@{}l@{}}#6\end{tabular}%
    \nodepart{three}\begin{tabular}{@{}l@{}}#7\end{tabular}};}

% \umlobjeto[opciones]{nombre-nodo}{x,y}{nombre:Clase}{slots}
\newcommand{\umlobjeto}[5][]{%
  \node[objeto,#1] (#2) at (#3) {\makebox[\linewidth][c]{\underline{\textbf{#4}}}%
    \nodepart{two}\begin{tabular}{@{}l@{}}#5\end{tabular}};}
 desde cada diagrama standalone.
\usepackage[T1]{fontenc}
\usepackage[utf8]{inputenc}
\usepackage{lmodern}
\usepackage{underscore}   % permite escribir identificadores con "_" sin escapar
\usepackage{xcolor}
\DeclareUnicodeCharacter{2192}{\ensuremath{\rightarrow}}
\DeclareUnicodeCharacter{2190}{\ensuremath{\leftarrow}}
\DeclareUnicodeCharacter{2194}{\ensuremath{\leftrightarrow}}
\DeclareUnicodeCharacter{2265}{\ensuremath{\geq}}
\DeclareUnicodeCharacter{2264}{\ensuremath{\leq}}
\DeclareUnicodeCharacter{2260}{\ensuremath{\neq}}
\DeclareUnicodeCharacter{00D7}{\ensuremath{\times}}
\DeclareUnicodeCharacter{21D2}{\ensuremath{\Rightarrow}}
\DeclareUnicodeCharacter{2026}{\ldots}
\usetikzlibrary{shapes.multipart,shapes.geometric,shapes.misc,arrows.meta,positioning,fit,backgrounds,calc}

% Paleta (sobria, imprime bien en B/N)
\definecolor{cPep}{HTML}{E3EEF9}     % agente / PEP
\definecolor{cPdp}{HTML}{FCEBD9}     % API / PDP
\definecolor{cDato}{HTML}{E4F2E2}    % datos / persistencia
\definecolor{cExt}{HTML}{EEE6F5}     % servicios externos
\definecolor{cUi}{HTML}{FFF6CC}      % frontend
\definecolor{cOff}{HTML}{EDEDED}     % offline
\definecolor{cBorde}{HTML}{3A3A3A}

\tikzset{
  % --- UML: clase de 3 compartimentos
  clase/.style={rectangle split, rectangle split parts=3, draw=cBorde, fill=white,
    font=\ttfamily\scriptsize, inner sep=2.5pt, text width=#1,
    rectangle split part align={center,left,left}},
  clase/.default=6cm,
  % --- UML: objeto (2 compartimentos, nombre subrayado)
  objeto/.style={rectangle split, rectangle split parts=2, draw=cBorde, fill=white,
    font=\ttfamily\scriptsize, inner sep=2.5pt, text width=#1,
    rectangle split part align={center,left}},
  objeto/.default=5cm,
  % --- relaciones UML
  herencia/.style={-{Triangle[open,length=7pt,width=8pt]}, draw=cBorde, thick},
  realiza/.style={-{Triangle[open,length=7pt,width=8pt]}, draw=cBorde, thick, dashed},
  compone/.style={{Diamond[length=9pt,width=6pt]}-{Stealth[length=5pt]}, draw=cBorde, thick},
  agrega/.style={{Diamond[open,length=9pt,width=6pt]}-{Stealth[length=5pt]}, draw=cBorde, thick},
  usa/.style={-{Stealth[length=5pt]}, draw=cBorde, dashed},
  asocia/.style={-{Stealth[length=5pt]}, draw=cBorde},
  enlace/.style={draw=cBorde, thick},
  rotulo/.style={font=\ttfamily\tiny, fill=white, inner sep=1pt},
  % --- componentes / flujo
  comp/.style={draw=cBorde, rounded corners=2pt, align=center, font=\small, inner sep=4pt, minimum height=9mm},
  paso/.style={draw=cBorde, rounded corners=3pt, align=center, font=\footnotesize, inner sep=3pt, text width=#1},
  paso/.default=3.4cm,
  decision/.style={draw=cBorde, diamond, aspect=2.2, align=center, font=\footnotesize, inner sep=1pt},
  almacen/.style={draw=cBorde, cylinder, shape border rotate=90, aspect=0.18, align=center,
    font=\footnotesize, minimum width=2.6cm, minimum height=1.2cm},
  zona/.style={draw=cBorde, dashed, rounded corners=4pt, inner sep=8pt},
  zonatit/.style={font=\bfseries\small, anchor=north west, inner sep=2pt},
  flujo/.style={-{Stealth[length=6pt]}, draw=cBorde, thick},
  flujofallo/.style={-{Stealth[length=6pt]}, draw=red!60!black, thick, dashed},
  proto/.style={font=\scriptsize\ttfamily, fill=white, inner sep=1pt, align=center},
}

% \umlclase[opciones]{nombre-nodo}{x,y}{estereotipo}{Nombre}{atributos}{operaciones}
% atributos/operaciones: líneas separadas con \\ ; usar \vacio si no hay.
\newcommand{\vacio}{\phantom{x}}
\newcommand{\umlclase}[7][]{%
  \node[clase,#1] (#2) at (#3) {\makebox[\linewidth][c]{%
    \begin{tabular}{@{}c@{}}\if\relax\detokenize{#4}\relax\else\textit{«#4»}\\\fi\textbf{\normalsize #5}\end{tabular}}%
    \nodepart{two}\begin{tabular}{@{}l@{}}#6\end{tabular}%
    \nodepart{three}\begin{tabular}{@{}l@{}}#7\end{tabular}};}

% \umlobjeto[opciones]{nombre-nodo}{x,y}{nombre:Clase}{slots}
\newcommand{\umlobjeto}[5][]{%
  \node[objeto,#1] (#2) at (#3) {\makebox[\linewidth][c]{\underline{\textbf{#4}}}%
    \nodepart{two}\begin{tabular}{@{}l@{}}#5\end{tabular}};}
 desde cada diagrama standalone.
\usepackage[T1]{fontenc}
\usepackage[utf8]{inputenc}
\usepackage{lmodern}
\usepackage{underscore}   % permite escribir identificadores con "_" sin escapar
\usepackage{xcolor}
\DeclareUnicodeCharacter{2192}{\ensuremath{\rightarrow}}
\DeclareUnicodeCharacter{2190}{\ensuremath{\leftarrow}}
\DeclareUnicodeCharacter{2194}{\ensuremath{\leftrightarrow}}
\DeclareUnicodeCharacter{2265}{\ensuremath{\geq}}
\DeclareUnicodeCharacter{2264}{\ensuremath{\leq}}
\DeclareUnicodeCharacter{2260}{\ensuremath{\neq}}
\DeclareUnicodeCharacter{00D7}{\ensuremath{\times}}
\DeclareUnicodeCharacter{21D2}{\ensuremath{\Rightarrow}}
\DeclareUnicodeCharacter{2026}{\ldots}
\usetikzlibrary{shapes.multipart,shapes.geometric,shapes.misc,arrows.meta,positioning,fit,backgrounds,calc}

% Paleta (sobria, imprime bien en B/N)
\definecolor{cPep}{HTML}{E3EEF9}     % agente / PEP
\definecolor{cPdp}{HTML}{FCEBD9}     % API / PDP
\definecolor{cDato}{HTML}{E4F2E2}    % datos / persistencia
\definecolor{cExt}{HTML}{EEE6F5}     % servicios externos
\definecolor{cUi}{HTML}{FFF6CC}      % frontend
\definecolor{cOff}{HTML}{EDEDED}     % offline
\definecolor{cBorde}{HTML}{3A3A3A}

\tikzset{
  % --- UML: clase de 3 compartimentos
  clase/.style={rectangle split, rectangle split parts=3, draw=cBorde, fill=white,
    font=\ttfamily\scriptsize, inner sep=2.5pt, text width=#1,
    rectangle split part align={center,left,left}},
  clase/.default=6cm,
  % --- UML: objeto (2 compartimentos, nombre subrayado)
  objeto/.style={rectangle split, rectangle split parts=2, draw=cBorde, fill=white,
    font=\ttfamily\scriptsize, inner sep=2.5pt, text width=#1,
    rectangle split part align={center,left}},
  objeto/.default=5cm,
  % --- relaciones UML
  herencia/.style={-{Triangle[open,length=7pt,width=8pt]}, draw=cBorde, thick},
  realiza/.style={-{Triangle[open,length=7pt,width=8pt]}, draw=cBorde, thick, dashed},
  compone/.style={{Diamond[length=9pt,width=6pt]}-{Stealth[length=5pt]}, draw=cBorde, thick},
  agrega/.style={{Diamond[open,length=9pt,width=6pt]}-{Stealth[length=5pt]}, draw=cBorde, thick},
  usa/.style={-{Stealth[length=5pt]}, draw=cBorde, dashed},
  asocia/.style={-{Stealth[length=5pt]}, draw=cBorde},
  enlace/.style={draw=cBorde, thick},
  rotulo/.style={font=\ttfamily\tiny, fill=white, inner sep=1pt},
  % --- componentes / flujo
  comp/.style={draw=cBorde, rounded corners=2pt, align=center, font=\small, inner sep=4pt, minimum height=9mm},
  paso/.style={draw=cBorde, rounded corners=3pt, align=center, font=\footnotesize, inner sep=3pt, text width=#1},
  paso/.default=3.4cm,
  decision/.style={draw=cBorde, diamond, aspect=2.2, align=center, font=\footnotesize, inner sep=1pt},
  almacen/.style={draw=cBorde, cylinder, shape border rotate=90, aspect=0.18, align=center,
    font=\footnotesize, minimum width=2.6cm, minimum height=1.2cm},
  zona/.style={draw=cBorde, dashed, rounded corners=4pt, inner sep=8pt},
  zonatit/.style={font=\bfseries\small, anchor=north west, inner sep=2pt},
  flujo/.style={-{Stealth[length=6pt]}, draw=cBorde, thick},
  flujofallo/.style={-{Stealth[length=6pt]}, draw=red!60!black, thick, dashed},
  proto/.style={font=\scriptsize\ttfamily, fill=white, inner sep=1pt, align=center},
}

% \umlclase[opciones]{nombre-nodo}{x,y}{estereotipo}{Nombre}{atributos}{operaciones}
% atributos/operaciones: líneas separadas con \\ ; usar \vacio si no hay.
\newcommand{\vacio}{\phantom{x}}
\newcommand{\umlclase}[7][]{%
  \node[clase,#1] (#2) at (#3) {\makebox[\linewidth][c]{%
    \begin{tabular}{@{}c@{}}\if\relax\detokenize{#4}\relax\else\textit{«#4»}\\\fi\textbf{\normalsize #5}\end{tabular}}%
    \nodepart{two}\begin{tabular}{@{}l@{}}#6\end{tabular}%
    \nodepart{three}\begin{tabular}{@{}l@{}}#7\end{tabular}};}

% \umlobjeto[opciones]{nombre-nodo}{x,y}{nombre:Clase}{slots}
\newcommand{\umlobjeto}[5][]{%
  \node[objeto,#1] (#2) at (#3) {\makebox[\linewidth][c]{\underline{\textbf{#4}}}%
    \nodepart{two}\begin{tabular}{@{}l@{}}#5\end{tabular}};}
